% !TEX root = ../main.tex
\chapter{变分法与 Euler 方程}
\label{ch:calcvar}

\noindent\emph{用到的前置工具：偏导数、链式法则与全微分（多元微积分部分）；分部积分（一元微积分）；无约束最优化的一阶条件（最优化部分）；差分方程与稳定性（前一章）。}
\medskip

到目前为止，本书谈的最优化几乎都是\emph{在有限个数}上做选择：选一个向量 $\vc x\in\RR^n$ 使目标最大。可经济学里有一大类问题，选的不是几个数，而是\emph{一整条随时间展开的路径}——一个人一辈子每个时点消费多少、一家企业每一年投资多少、一个经济体如何随时间积累资本。这时决策变量是一个\emph{函数} $x(\cdot)$，目标是把它代进某个积分后得到的一个数。在所有可行路径里挑出使这个积分取极值的那一条，就是\textbf{变分法}（calculus of variations）要解决的问题。本章先说清这类问题长什么样、为什么普通求导用不上，再把它的核心工具——Euler\,–\,Lagrange 方程——从头推出来，并交代端点自由时补上的横截条件。学完会看到：连续时间宏观里反复出现的``消费 Euler 方程''，正是本章方程的一个特例；而下一章的最优控制（Pontryagin），则是把这里的思路再推广一层。

\section{从``选一个数''到``选一条路径''}
\label{sec:calcvar-1}

先看一个具体场景。一个人计划在 $[0,T]$ 这段时间里安排消费 $c(t)$，瞬时效用为 $u(c)$、主观贴现率为 $\rho$，他想最大化一生的贴现效用
\[
    \int_0^T e^{-\rho t}\,u\bigl(c(t)\bigr)\,\d t .
\]
这里要选的不是一个消费水平，而是\emph{每个时点的消费}——也就是整条函数 $c:[0,T]\to\RR$。目标也不再是一个普通函数 $f(c)$，而是``把一整条曲线喂进去、吐出一个数''的东西。我们把这种``函数的函数''叫\textbf{泛函}。

\defn{泛函与最简泛函问题}{
设 $L(t,x,v)$ 是一个三元连续可微函数（称为\textbf{被积函数}或 Lagrange 量），端点 $a<b$ 与端点值 $x_a,x_b$ 给定。对每条连续可微路径 $x:[a,b]\to\RR$，定义\textbf{泛函}
\[
    J[x]\;=\;\int_a^b L\bigl(t,\,x(t),\,\dot x(t)\bigr)\,\d t ,
\]
其中 $\dot x=\d x/\d t$。\emph{最简变分问题}是在满足端点条件 $x(a)=x_a$、$x(b)=x_b$ 的所有路径中，求使 $J[x]$ 取极值（极小或极大）的那条 $x^*$。
}

记号上要分清两层括号：$J[x]$ 用方括号，强调输入是\emph{整条函数}；而 $L(t,x,\dot x)$ 用圆括号，它是普通的多元函数，只不过我们习惯把它的三个槽位分别叫``时间''$t$、``状态''$x$、``速度''$\dot x$。求 $\partial L/\partial x$ 与 $\partial L/\partial\dot x$ 时，就把 $L$ 当成以 $t,x,\dot x$ 为\emph{三个独立变量}的普通函数来偏导——这一点初学时最容易绊住，务必记牢。

\state{为什么不能直接用普通求导}{
普通最优化是``目标对一个数求导、令其为零''。可这里的自变量是\emph{一条函数}，没有``一个数''可供求导。变分法的整套技巧，本质上是把``在函数空间里取极值''这件无从下手的事，化归成``沿每一个允许的扰动方向取极值''——而每个方向上又退回成熟悉的一元求导。下一节就把这条化归落实成一个可计算的微分方程。
}

这类问题在经济学里的三大主场，正好对应三个名字：\textbf{最优增长}（Ramsey 模型的连续时间版，选资本积累路径）、\textbf{最优消费}（选一生的消费路径）、以及\textbf{带调整成本的投资}（选投资路径，使建造成本与调整成本的现值最小）。它们的共同点都是``在一整条时间路径上做选择'',这正是变分法的用武之地。

\section{Euler\,–\,Lagrange 方程}
\label{sec:calcvar-2}

\subsection{取变分：把函数极值化归成一元求导}

设 $x^*$ 是问题的解。变分法的核心想法是：既然 $x^*$ 是``最好''的路径，那么从它出发、朝\emph{任何允许的方向}做一点点偏移，目标都不该变得更好。怎么描述``一个方向''？取一条任意的、在两端为零的扰动函数 $\eta(t)$，
\[
    \eta(a)=0,\qquad \eta(b)=0 ,
\]
两端为零是为了让偏移后的路径仍满足同样的端点条件。再用一个小标量 $\ve$ 控制偏移幅度，构造\emph{一族}邻近路径
\[
    x_\ve(t)\;=\;x^*(t)+\ve\,\eta(t) .
\]
当 $\ve=0$ 时回到最优路径，$\ve\neq0$ 时是它的一个扰动（图~\ref{fig:calcvar-1}）。把这一族路径代进泛函，得到一个\emph{只依赖于一个数 $\ve$}的普通函数：
\[
    \Phi(\ve)\;=\;J[x^*+\ve\eta]\;=\;\int_a^b L\bigl(t,\,x^*+\ve\eta,\,\dot x^*+\ve\dot\eta\bigr)\,\d t .
\]

\begin{figure}[h]
\centering
\begin{tikzpicture}[scale=1.15,>=Stealth]
    % 坐标轴
    \draw[->] (-0.2,0) -- (5.0,0) node[right] {$t$};
    \draw[->] (0,-0.2) -- (0,5.2) node[above] {$x$};
    % 端点刻度
    \node[anchor=east] at (-0.05,1) {$x_a$};
    \node[anchor=east] at (-0.05,4.12) {$x_b$};
    \draw[dashed,gray!60] (0,1) -- (1.0,1);
    \draw[dashed,gray!60] (0,4.12) -- (4,4.12);
    \node[anchor=north] at (1.0,-0.05) {$a$};
    \node[anchor=north] at (4.0,-0.05) {$b$};
    \draw[dashed,gray!60] (1.0,0) -- (1.0,1);
    \draw[dashed,gray!60] (4.0,0) -- (4.0,4.12);
    % 邻近变分路径（虚线）：x* ± 0.6 sin(pi (t-a)/3)，端点与 x* 重合
    \draw[densely dashed,red!55!black,domain=1.0:4.0,smooth,samples=80,variable=\t]
        plot ({\t},{1+1.5*(\t-1)-0.18*(\t-1)*(\t-1)+0.6*sin(deg(pi*(\t-1)/3))});
    \draw[densely dashed,red!55!black,domain=1.0:4.0,smooth,samples=80,variable=\t]
        plot ({\t},{1+1.5*(\t-1)-0.18*(\t-1)*(\t-1)-0.6*sin(deg(pi*(\t-1)/3))});
    % 最优路径（实线）
    \draw[very thick,blue!60!black,domain=1.0:4.0,smooth,samples=80,variable=\t]
        plot ({\t},{1+1.5*(\t-1)-0.18*(\t-1)*(\t-1)});
    % 固定端点
    \fill (1.0,1) circle (1.8pt);
    \fill (4.0,4.12) circle (1.8pt);
    % 标签：最优路径，停在曲线下方空白，细引线指向曲线
    \node[blue!60!black,anchor=north] at (2.7,2.05) {$x^*(t)$};
    \draw[blue!60!black,->] (2.7,2.1) -- (2.55,2.74);
    % 标签：上方变分路径
    \node[red!55!black,anchor=south west] at (1.5,3.7) {$x^*(t)+\varepsilon\eta(t)$};
    \draw[red!55!black,->] (2.05,3.72) -- (2.2,3.55);
    % 标签：下方变分路径
    \node[red!55!black,anchor=north west] at (2.1,1.5) {$x^*(t)-\varepsilon\eta(t)$};
    \draw[red!55!black,->] (2.55,1.6) -- (2.6,2.0);
\end{tikzpicture}
\caption{最优路径 $x^*(t)$（实线）与两条端点重合的邻近变分路径 $x^*\pm\ve\eta$（虚线）。扰动 $\eta$ 在两端 $a,b$ 为零，故所有路径共用同样的端点。}
\label{fig:calcvar-1}
\end{figure}

现在 $\Phi$ 是一个普通的一元函数，$x^*$ 最优意味着 $\ve=0$ 是 $\Phi$ 的极值点。于是它必须满足熟悉的一元一阶条件
\[
    \Phi'(0)=0 .
\]
几何上，把 $\Phi$ 画成 $\ve$ 的曲线（图~\ref{fig:calcvar-2}），最优性要求它在 $\ve=0$ 处\emph{驻定}——切线水平。这一步是整套方法的支点：函数空间里``取极值''这件抽象的事，被压成了``一族一元曲线全都在原点驻定''。

\begin{figure}[h]
\centering
\begin{tikzpicture}[scale=1.15,>=Stealth]
    % 坐标轴：横轴 eps，纵轴 Phi(eps)=J[x*+eps eta]
    \draw[->] (-2.6,0) -- (2.7,0) node[right] {$\varepsilon$};
    \draw[->] (0,-0.3) -- (0,3.4) node[above] {$\Phi(\varepsilon)=J[x^*+\varepsilon\eta]$};
    % 抛物线 Phi(eps)=1.5+0.35 eps^2，最小在 eps=0
    \draw[very thick,blue!60!black,domain=-2.35:2.35,smooth,samples=80,variable=\e]
        plot ({\e},{1.5+0.35*\e*\e});
    % eps=0 处的水平切线（驻点）
    \draw[red!60!black,thick] (-1.5,1.5) -- (1.5,1.5);
    % 驻点
    \fill (0,1.5) circle (1.8pt);
    \draw[dashed,gray!60] (0,1.5) -- (0,0) ;
    \node[anchor=north] at (0,-0.05) {$0$};
    % 标签：水平切线 = 一阶变分为零
    \node[red!60!black,anchor=south] at (0,1.55) {$\Phi'(0)=0$};
\end{tikzpicture}
\caption{沿任一扰动方向 $\eta$，泛函的值 $\Phi(\ve)=J[x^*+\ve\eta]$ 是 $\ve$ 的普通函数。最优路径使它在 $\ve=0$ 处驻定（切线水平），即一阶变分为零。}
\label{fig:calcvar-2}
\end{figure}

\subsection{算出 $\Phi'(0)$ 并分部积分}

把 $\Phi$ 对 $\ve$ 求导。积分号下对 $\ve$ 求导，再用链式法则（注意 $L$ 的第二槽 $x$ 处增量是 $\ve\eta$、第三槽 $\dot x$ 处增量是 $\ve\dot\eta$）：
\[
    \Phi'(\ve)=\int_a^b\br{\pd{L}{x}\bigl(t,x_\ve,\dot x_\ve\bigr)\,\eta(t)
    +\pd{L}{\dot x}\bigl(t,x_\ve,\dot x_\ve\bigr)\,\dot\eta(t)}\d t .
\]
令 $\ve=0$（此时 $x_\ve=x^*$、$\dot x_\ve=\dot x^*$），并把在 $x^*$ 处取值的偏导简记为 $L_x,L_{\dot x}$，一阶条件 $\Phi'(0)=0$ 化为
\begin{equation}
    \int_a^b\br{L_x\,\eta+L_{\dot x}\,\dot\eta}\,\d t=0 ,
    \qquad\text{对一切满足 }\eta(a)=\eta(b)=0\text{ 的 }\eta .
    \label{eq:calcvar-firstvar}
\end{equation}
左边这个量叫泛函的\textbf{一阶变分}，常记作 $\delta J$。现在的麻烦是：式~\eqref{eq:calcvar-firstvar} 里同时出现了 $\eta$ 和它的导数 $\dot\eta$，无法直接把 $\eta$ 当作``任意''提出来。出路是对含 $\dot\eta$ 的那一项做\emph{分部积分}，把导数从 $\eta$ 身上``搬''到 $L_{\dot x}$ 身上：
\[
    \int_a^b L_{\dot x}\,\dot\eta\,\d t
    =\underbrace{\br{L_{\dot x}\,\eta}_a^b}_{=\,0}
    -\int_a^b\dd{}{t}\pr{L_{\dot x}}\,\eta\,\d t .
\]
边界项 $\br{L_{\dot x}\eta}_a^b$ 因 $\eta(a)=\eta(b)=0$ 而消失——这正是``两端为零''的扰动派上的用场。代回式~\eqref{eq:calcvar-firstvar}，把两项合并、提出公因子 $\eta$：
\begin{equation}
    \int_a^b\br{\,L_x-\dd{}{t}\pr{L_{\dot x}}\,}\eta(t)\,\d t=0 ,
    \qquad\text{对一切两端为零的 }\eta .
    \label{eq:calcvar-needlemma}
\end{equation}

\subsection{变分基本引理：从``积分恒为零''到``被积式恒为零''}

式~\eqref{eq:calcvar-needlemma} 说的是``某个固定函数乘以任意 $\eta$、积分总为零''。直觉上，能对\emph{所有}扰动都正交的函数只能是零函数。这一步要靠一条小引理把它说严密。

\lem{变分基本引理}{
设 $h:[a,b]\to\RR$ 连续。若对一切满足 $\eta(a)=\eta(b)=0$ 的连续可微函数 $\eta$ 都有 $\int_a^b h(t)\,\eta(t)\,\d t=0$，则 $h(t)\equiv 0$ 于 $[a,b]$。
}
\pf{
反证。设某点 $t_0\in(a,b)$ 处 $h(t_0)\neq0$，不妨 $h(t_0)>0$。由连续性，存在 $t_0$ 的一个小区间 $(c,d)\subseteq(a,b)$，在其上 $h>0$。专门挑一个``只在这小区间上凸起、别处为零''的扰动：令
\[
    \eta(t)=
    \begin{cases}
        (t-c)^2(d-t)^2, & t\in(c,d),\\[2pt]
        0, & \text{其它},
    \end{cases}
\]
它连续可微、两端为零、且在 $(c,d)$ 内严格为正。于是 $\int_a^b h\,\eta\,\d t=\int_c^d h\,\eta\,\d t>0$（正乘正的积分为正），与假设矛盾。故 $h$ 处处为零。
}

\rmkb[这条引理的``经济''读法]{
把 $\eta$ 想成``你可以朝任意方向、任意局部地推一把''的检验扰动。若被积式 $h$ 在某处不为零，你总能精准地``只在那一处推一把''而把积分顶离零——这正是上面那个局部凸起函数干的事。反过来说，唯一能扛住一切局部扰动的，只有处处为零。变分法里``任意扰动''的威力，全在于此。
}

把引理用到式~\eqref{eq:calcvar-needlemma}（其中 $h=L_x-\frac{\d}{\d t}L_{\dot x}$）上，立刻得到本章的主结果。

\thm{Euler\,–\,Lagrange 方程}{
若连续可微路径 $x^*$ 在满足端点条件 $x(a)=x_a,\,x(b)=x_b$ 的路径中使 $J[x]=\int_a^b L(t,x,\dot x)\,\d t$ 取极值，则它在 $[a,b]$ 上满足
\[
    \boxed{\ \dd{}{t}\pr{\pd{L}{\dot x}}=\pd{L}{x}\ } .
\]
}

这是一个关于未知函数 $x(t)$ 的二阶常微分方程：把全导数 $\frac{\d}{\d t}$ 展开（$L_{\dot x}$ 通过 $t,x,\dot x$ 依赖 $t$），会出现 $\ddot x$。它配上两个端点条件，原则上就把最优路径定了下来。

\rmkb[这是必要条件，不是充分条件]{
Euler\,–\,Lagrange 方程来自一阶条件 $\Phi'(0)=0$，因此只是``驻定''的必要条件——它筛出的是极小、极大与拐点的候选。要断定确实是极小（或极大），需要二阶条件（沿任意 $\eta$ 有 $\Phi''(0)\ge0$）：其必要形式是 Legendre 条件 $L_{\dot x\dot x}\ge0$（求极小时），强化为 $L_{\dot x\dot x}>0$ 并附 Jacobi（共轭点）条件后才充分。本书按使用者的需要，重点放在用 E\,–\,L 方程\emph{解出}候选路径上；经济学问题里目标函数往往天然凹（效用凹、成本凸），驻定路径即为所求，二阶条件多半自动满足。
}

\subsection{两个能省力的特例}

E\,–\,L 方程在两类常见结构下可以``降一阶'',值得单独记住。

\state{两个守恒型特例}{
\textbf{(i) $L$ 不显含 $x$}（只含 $t,\dot x$）。则 $\partial L/\partial x=0$，E\,–\,L 退化为 $\frac{\d}{\d t}(\partial L/\partial\dot x)=0$，即
\[
    \pd{L}{\dot x}=\text{常数}.
\]
这个守恒量是``广义动量''。\quad
\textbf{(ii) $L$ 不显含 $t$}（只含 $x,\dot x$，称\emph{自治}）。则有 Beltrami 恒等式
\[
    L-\dot x\,\pd{L}{\dot x}=\text{常数}.
\]
两种情形都把二阶方程降成一阶，往往能直接积分出来。
}

特例 (ii) 的来历值得一推，因为它示范了``沿最优路径某个量守恒''这一在动态经济学里反复出现的母题。沿任一路径，按链式法则
\[
    \dd{}{t}\pr{L-\dot x\,L_{\dot x}}
    =\underbrace{\bigl(L_t+L_x\dot x+L_{\dot x}\ddot x\bigr)}_{\d L/\d t}
    -\ddot x\,L_{\dot x}-\dot x\,\dd{}{t}L_{\dot x} .
\]
自治时 $L_t=0$；再把 E\,–\,L 方程 $\frac{\d}{\d t}L_{\dot x}=L_x$ 代入末项，得
\[
    \dd{}{t}\pr{L-\dot x\,L_{\dot x}}
    =L_x\dot x+L_{\dot x}\ddot x-\ddot x\,L_{\dot x}-\dot x\,L_x=0 .
\]
故 $L-\dot x\,L_{\dot x}$ 沿最优路径恒定。

\rmk{本节守恒的量是 $L-\dot x\,L_{\dot x}$；把它取负号得到 $\dot x\,L_{\dot x}-L$，正是后文将出现的 \emph{Hamilton 量} $H$ 的雏形（在物理里即能量）。变分法里``时间不显含 $\Rightarrow$ 某量守恒''的现象，在最优控制与物理里都对应``Hamilton 量守恒''。}

\section{横截条件：端点自由时补上的边界}
\label{sec:calcvar-3}

上一节假设两端点 $x_a,x_b$ 都给定，所以扰动被迫两端为零。可经济学里末端往往是\emph{自由}的：一个人到了寿命终点 $T$，留多少遗产 $x(T)$ 并不外生给定，而要由最优化\emph{内生决定}。这时 $\eta(b)$ 不必为零，分部积分里那个边界项就不会自动消失，反而会给我们补上一个缺失的边界条件——\textbf{横截条件}（transversality condition）。

把端点自由的情形重做一遍：左端 $x(a)=x_a$ 仍固定（$\eta(a)=0$），右端 $x(b)$ 自由（$\eta(b)$ 任意）。一阶变分经分部积分后是
\[
    \Phi'(0)=\int_a^b\br{L_x-\dd{}{t}L_{\dot x}}\eta\,\d t
    +\br{L_{\dot x}\,\eta}_a^b=0 .
\]
边界项现在只剩右端 $L_{\dot x}(b)\,\eta(b)$（左端因 $\eta(a)=0$ 仍为零）。要让\emph{所有}允许的 $\eta$ 都使 $\Phi'(0)=0$，先取那些恰好 $\eta(b)=0$ 的扰动——它们逼出内部的 E\,–\,L 方程照旧成立，于是积分项归零；剩下的就要求边界项对\emph{任意} $\eta(b)$ 都为零。这只能是
\[
    \boxed{\ \pd{L}{\dot x}\Big|_{t=b}=0\ } .
\]

\state{横截条件的来历与含义}{
横截条件不是额外强加的，它是\emph{同一个}一阶变分在``端点不再被钉死''时自然吐出的边界条件——端点一旦获得自由，相应的边界项就必须独立为零。它恰好补上了因放开一个端点条件而``少掉''的那一个定解条件，使二阶 E\,–\,L 方程仍有唯一解。在不同设定下它换不同的面孔：末端值自由时是 $L_{\dot x}(b)=0$；带贴现、末端要求 $x(T)\ge0$（如不许负遗产）时，化为带互补松弛的 $e^{-\rho T}\lambda(T)\,x(T)=0$ 形式——也就是宏观里那条``资本的影子价值现值在终点归零''的著名条件。
}

\rmk{无穷期问题（$T\to\infty$）里横截条件取极限形式 $\lim_{t\to\infty}e^{-\rho t}\lambda(t)\,x(t)=0$，用来排除``庞氏式''地无限累积或挥霍资本的发散路径。它在连续时间增长与最优控制里是判定路径是否最优的关键一环，下一章会再正式登场。}

\section{应用：最优消费路径}
\label{sec:calcvar-4}

现在用 E\,–\,L 方程把本章开头那个消费问题解到底，看它如何吐出宏观里天天用的``消费 Euler 方程''。

\ex[贴现效用下的最优消费增长]{
个体在 $[0,T]$ 内持有资本 $k(t)$，以利率 $r$ 取得资本收益，预算约束（资本运动方程）为
\[
    \dot k=r\,k-c ,\qquad k(0)=k_0,\ k(T)=k_T\ \text{给定}.
\]
他最大化贴现效用 $\displaystyle\int_0^T e^{-\rho t}u(c)\,\d t$，其中 $u$ 严格增、严格凹。求最优消费路径满足的方程；再取等弹性效用 $u(c)=\dfrac{c^{1-\sigma}}{1-\sigma}$（$\sigma>0$，跨期替代弹性的倒数）算出消费的增长率。
}
\sol{
约束 $\dot k=rk-c$ 把消费表成 $c=rk-\dot k$。代入目标，问题就化成了一个以 $k$ 为路径、最简形式的变分问题，被积函数为
\[
    L(t,k,\dot k)=e^{-\rho t}\,u\!\bigl(rk-\dot k\bigr).
\]
算两个偏导（对 $L$ 的第二槽 $k$ 与第三槽 $\dot k$，记 $u'=u'(c)$）：
\[
    \pd{L}{k}=e^{-\rho t}u'(c)\cdot r ,
    \qquad
    \pd{L}{\dot k}=e^{-\rho t}u'(c)\cdot(-1)=-e^{-\rho t}u'(c).
\]
代进 E\,–\,L 方程 $\dfrac{\d}{\d t}\!\pr{\partial L/\partial\dot k}=\partial L/\partial k$：
\[
    \dd{}{t}\br{-e^{-\rho t}u'(c)}=r\,e^{-\rho t}u'(c).
\]
左边求导（乘积法则，$u'(c)$ 再经 $c$ 依赖 $t$）：
\[
    \rho\,e^{-\rho t}u'(c)-e^{-\rho t}u''(c)\,\dot c=r\,e^{-\rho t}u'(c).
\]
两边除以 $e^{-\rho t}$ 并整理，得\textbf{消费 Euler 方程}
\[
    -\frac{u''(c)}{u'(c)}\,\dot c=r-\rho .
\]
它的读法很干净：左边是``消费每变动一点带来的边际效用相对变化'',右边是``市场回报 $r$ 与不耐心 $\rho$ 之差''。当 $r>\rho$（市场回报跑赢贴现），个体愿意推迟消费、让消费\emph{随时间上升}；$r<\rho$ 则相反；$r=\rho$ 时消费平滑为常数。

取等弹性效用 $u(c)=c^{1-\sigma}/(1-\sigma)$，则 $u'=c^{-\sigma}$、$u''=-\sigma c^{-\sigma-1}$，于是 $-u''/u'=\sigma/c$，左边成 $\sigma\,\dot c/c$。代入得消费增长率
\[
    \boxed{\ \frac{\dot c}{c}=\frac{r-\rho}{\sigma}\ } .
\]
增长率为常数，故消费沿指数路径 $c(t)=c_0\,e^{(r-\rho)t/\sigma}$ 演化，三种情形如图~\ref{fig:calcvar-3}。$\sigma$ 越大（越不愿在时间上替代消费），同样的 $r-\rho$ 推动的增长越平缓。
}

\begin{figure}[h]
\centering
\begin{tikzpicture}[xscale=1.25,yscale=1.05,>=Stealth]
    % 坐标轴：横轴 t，纵轴 c(t)
    \draw[->] (-0.2,0) -- (4.7,0) node[right] {$t$};
    \draw[->] (0,-0.2) -- (0,5.0) node[above] {$c(t)$};
    % 三条路径，c0=1.6，g=0.25 / 0 / -0.25
    \draw[very thick,red!70!black,domain=0:4.0,smooth,samples=80,variable=\t]
        plot ({\t},{1.6*exp(0.25*\t)});
    \draw[very thick,blue!60!black,domain=0:4.0,smooth,samples=80,variable=\t]
        plot ({\t},{1.6});
    \draw[very thick,green!50!black,domain=0:4.0,smooth,samples=80,variable=\t]
        plot ({\t},{1.6*exp(-0.25*\t)});
    % 共同起点
    \fill (0,1.6) circle (1.8pt);
    \node[anchor=east] at (-0.05,1.6) {$c_0$};
    % 曲线标签，停在各曲线右端空白
    \node[red!70!black,anchor=west] at (4.05,4.35) {$r>\rho$};
    \node[blue!60!black,anchor=west] at (4.05,1.6) {$r=\rho$};
    \node[green!50!black,anchor=west] at (4.05,0.59) {$r<\rho$};
    % Euler 方程标注，停在左上空白
    \node[anchor=west] at (0.25,4.55) {$\dfrac{\dot c}{c}=\dfrac{r-\rho}{\sigma}$};
\end{tikzpicture}
\caption{等弹性效用下的最优消费路径 $c(t)=c_0\,e^{(r-\rho)t/\sigma}$：$r>\rho$ 时上升、$r=\rho$ 时平直、$r<\rho$ 时下降。斜率由 $r-\rho$ 定方向、由跨期替代参数 $\sigma$ 定陡缓。}
\label{fig:calcvar-3}
\end{figure}

\rmkb[本例用了哪些前置工具]{
注意整条求解链：用约束把控制 $c$ 消成只剩状态路径 $k$（化归成最简问题）、对三槽函数 $L$ 取偏导（多元微积分）、用乘积与链式法则展开全导数、最后解出一个一阶常微分方程得增长率（前一章的差分方程在连续时间的对应物）。这正是本书反复强调的：一个``高级''结果，拆开都是几件基础工具的组合。
}

\section{它通向何处}
\label{sec:calcvar-5}

变分法是连续时间动态最优化的\emph{第一块基石}，它的几处主要去处把本章接回更大的版图：

\begin{itemize}
    \item \textbf{最优控制与 Pontryagin 最大值原理}（下一章）。本章把控制变量 $c$ ``塞回''约束、化成纯状态的变分问题；可一旦约束复杂（控制有上下界、状态受不等式约束），这条``代入消元''的路就走不通了。最优控制把状态 $k$ 与控制 $c$ 分开处理，引入\emph{协态变量}（影子价格）$\lambda$ 与 \emph{Hamilton 量} $\mathcal H=u(c)+\lambda(rk-c)$，最优性条件 $\partial\mathcal H/\partial c=0$、$\dot\lambda=\rho\lambda-\partial\mathcal H/\partial k$ 正是 E\,–\,L 方程的推广——把 $\partial\mathcal H/\partial c=0$ 与协态方程联立消去 $\lambda$，又会回到本章的消费 Euler 方程。变分法是它的特例：``无约束''或``约束可代入''的那一类。
    \item \textbf{连续时间增长与 Hamilton 量}。Ramsey\,–\,Cass\,–\,Koopmans 最优增长模型的连续时间版，骨架就是本章这套：被积函数是贴现效用、约束是资本运动方程、解出的就是消费与资本的 Euler 方程加横截条件。本章自治情形下守恒的 $L-\dot x\,L_{\dot x}$，正是 Hamilton 量的前身。
    \item \textbf{消费 Euler 方程，连续与离散}。本章解出的 $\dot c/c=(r-\rho)/\sigma$ 是宏观与资产定价里最常被引用的关系之一。它的离散时间表亲 $u'(c_t)=\beta(1+r)\,\E{u'(c_{t+1})}$ 来自动态规划的 Bellman 方程（见后文``动态规划''一章），两者讲的是同一件事——\emph{边际效用沿最优路径如何随时间移动}——只是一个用微分方程、一个用差分方程表达。计量上，把后者的随机版本作为矩条件去估参数，就是宏观结构估计里经典的 Euler 方程 GMM（见渐近统计部分的极值估计框架）。
\end{itemize}

一条贯穿的主线于是浮现：\emph{在整条时间路径上做选择}，离散版是动态规划、连续版是变分法与最优控制，而它们最终都落到同一句话上——沿最优路径，相邻时点的边际权衡必须恰好相等。本章给了这句话最直接的微分形式。
